% \vspace{-0.2cm}
\section{Conclusion}
\label{sec:conclusion}
This paper performs the most extensive systematic study of the ML-based Android malware detection literature with empirical and quantitative analysis.
We identify challenges (\textit{i.e.,} unfair comparisons, unrealistic evaluations, and unclear computational costs) that hinder the systematization in this field.
In response, we design a general-purpose framework for developing ML-based detection approaches and evaluating their effectiveness, robustness, and efficiency.
By experimentally comparing 12 representative approaches, our study paints a holistic view of the state of ML-based Android malware detection and puts forth recommendations to guide future research.
For instance, we find existing detectors are still vulnerable to malware evolution and adversarial attacks and in future work, we should focus on incorporating more \textit{malware semantics} to design more practical detectors.
% Committed to the research community's growth, all the artifacts (code, data, and logs) are available at \url{https://anonymous.4open.science/r/RkRyb2lk}.
% Committed to the research community's growth, all the artifacts (code, data, and logs) will be released upon acceptance.

\section*{Acknowledgments}
We thank Bonan Ruan, Jiawei Li, and Chuqi Zhang for their assistance and the anonymous reviewers for their valuable comments.
This research is supported by the National Research Foundation, Singapore, through the National Cybersecurity R\&D Lab at the National University of Singapore under its National Cybersecurity R\&D Programme (Award No. NCR25-NCL P3-0001). 
Any opinions, findings and conclusions or recommendations expressed in this material are those of the authors(s) and do not reflect the views of National Research Foundation, Singapore.